% Part: computability
% Chapter: recursive-functions
% Section: notations-pr-functions

\documentclass[../../../include/open-logic-section]{subfiles}

\begin{document}

% SEGMENT OLP-0214-S01

\olfileid{cmp}{rec}{not}
\olsection{முதன்மை மீள்வரையறைக் குறிமுறைகள்}

முதன்மை மீள்வரையறைச் சார்புகளின் துல்லியமான தொகுத்தறிதல் விளக்கத்தைக்
கொண்டிருப்பதன் ஒரு நன்மை, அவற்றை முறையாக விவரிக்க முடிவதாகும்.
எடுத்துக்காட்டாக, அத்தகைய ஒவ்வொரு சார்புக்கும் பின்வருமாறு ஒரு
“குறிமுறையை” ஒதுக்கலாம். சுழியம், அடுத்தெண், வீழ்த்தல்கள் ஆகியவற்றுக்கு
முறையே $\Zero$, $\Succ$ மற்றும் $\Proj{n}{i}$ குறியீடுகளைப் பயன்படுத்துக.
இப்போது $k$-இடச் சார்பு~$f$ மற்றும் $n$-இடச் சார்புகள் $g_0$,
\dots,~$g_{k-1}$ ஆகியவற்றிலிருந்து சேர்ப்பால் $h$
வரையறுக்கப்படுகிறது என்றும், பிந்தைய சார்புகளுக்கு $F$, $G_0$,
\dots,~$G_{k-1}$ என்ற குறிமுறைகளை ஒதுக்கியுள்ளோம் என்றும் கொள்க.
அப்போது புதிய குறியீடு $\fn{Comp}_{k,n}$ ஐப் பயன்படுத்தி,
சார்பு $h$ ஐ $\fn{Comp}_{k,n}[F,G_0,\dots,G_{k-1}]$ எனக் குறிக்கலாம்.

முதன்மை மீள்வரையறையால் வரையறுக்கப்படும் சார்புகளுக்கு ஒத்த
குறிமுறைகளைப் பயன்படுத்தலாம். $(k+1)$-இடச் சார்பு~$h$ ஆனது
$k$-இடச் சார்பு~$f$ மற்றும் $(k+2)$-இடச் சார்பு~$g$ இலிருந்து
முதன்மை மீள்வரையறையால் வரையறுக்கப்படுகிறது என்றும், $f$ மற்றும்
~$g$ க்கு ஒதுக்கப்பட்ட குறிமுறைகள் முறையே $F$ மற்றும்~$G$
என்றும் கொள்க. அப்போது~$h$ க்கு ஒதுக்கப்படும் குறிமுறை
$\fn{Rec}_k[F,G]$.

கூட்டல் சார்பு முதன்மை மீள்வரையறையால் பின்வருமாறு வரையறுக்கப்படுவதை
நினைவுகூர்க:
\begin{align*}
  \Add(x_0, 0) & = \Proj{1}{0}(x_0) = x_0\\
  \Add(x_0, y+1) & = \Succ(\Proj{3}{2}(x_0, y, \Add(x_0, y))) = \Add(x_0, y) +1
\end{align*}
இங்கே~$f$ இன் பணியை $\Proj{1}{0}$ வகிக்கிறது;~$g$ இன் பணியை
$\Succ(\Proj{3}{2}(x_0, y, z))$ வகிக்கிறது. பிந்தையது
$1$-இடச் சார்பு~$\Succ$ மற்றும் $3$-இடச் சார்பு~$\Proj{3}{2}$
ஆகியவற்றிலிருந்து சேர்ப்பால் வரையறுக்கப்படும் சார்பின் விளைவு என்பதால்,
அதற்கு $\fn{Comp}_{1,3}[\Succ,\Proj{3}{2}]$ என்ற குறிமுறை
ஒதுக்கப்படுகிறது. இவ்வமைப்பில், கூட்டல் சார்பைப் பின்வருமாறு குறிக்கலாம்:
\[
\fn{Rec}_1[\Proj{1}{0},\fn{Comp}_{1,3}[\Succ,\Proj{3}{2}]].
\]
இக்குறிமுறைகள் சிலவேளைகளில் பயனளிக்கின்றன; எடுத்துக்காட்டாக,
முதன்மை மீள்வரையறைச் சார்புகளை எண்ணிடும்போது.

\begin{prob}
  ~$\Mult$ க்கான முழுமையான முதன்மை மீள்வரையறைக் குறிமுறையைத் தருக.
\end{prob}

\end{document}
